%! Semi-log Plots — Unit 2, Lesson 2.15 — Homework (Answer Key)
\documentclass[10pt]{article}
\usepackage{precalculus-article}
\usepackage{precalculus-key}   % pulls in -boxes; do NOT also load -boxes
\newcommand{\TallMath}[1]{$\displaystyle #1\rule[-1.4em]{0pt}{3.2em}$}

\begin{document}

\pageheader{Unit 2, Lesson 2.15}{Homework --- Answer Key}
\namedateperiod

\vspace{0.1in}

\begin{enumerate}[leftmargin=1.8em, itemsep=14pt]

\item \textbf{Conceptual.}  Explain in your own words why exponential data appears linear
      on a semi-log plot.

\ansline{When you apply $\log$ to both sides of $y = ab^x$, the exponent $x$ comes out of
the logarithm (power rule), giving $\log(y) = x\log(b) + \log(a)$. This is a linear function
of $x$, so plotting $\log(y)$ against $x$ on a semi-log graph produces a straight line. The
multiplicative growth of $y$ is converted to additive growth of $\log(y)$.}

\item Given the exponential model $y = 2 \cdot 5^x$, write the linearized form using
      $\log_{10}$.  Identify the slope and $y$-intercept.

\vspace{4pt}
Linearized form: $\log_{10}(y) = $ \ans{$(\log_{10} 5)\cdot x + \log_{10} 2 \approx 0.699x + 0.301$}

\vspace{4pt}
Slope $= $ \ans{$\log_{10}(5) \approx 0.699$} \qquad $y$-intercept $= $ \ans{$\log_{10}(2) \approx 0.301$}

\item A semi-log plot with slope $0.477$ and $y$-intercept $1.0$.  Write the exponential model.

\vspace{4pt}
$b = 10^{0.477} \approx $ \ans{$3$} \qquad $a = 10^{1.0} = $ \ans{$10$}

\vspace{4pt}
Exponential model: $y = $ \ans{$10 \cdot 3^x$}

\item The table below gives data for a population $P$ over time $t$.

\vspace{4pt}
\begin{center}
\renewcommand{\arraystretch}{2.0}
\begin{tabular}{|c|c|c|c|c|}
\hline
\rowcolor{sky}
$t$ & $0$ & $2$ & $4$ & $6$ \\
\hline
$P$ & $500$ & $2000$ & $8000$ & $32000$ \\
\hline
$\log_{10}(P)$ & \ans{2.699} & \ans{3.301} & \ans{3.903} & \ans{4.505} \\
\hline
\end{tabular}
\end{center}

\begin{enumerate}[leftmargin=1.5em, itemsep=6pt, label=\alph*.]
  \item (Completed in table above.)
  \item Slope $= $ \ans{$(3.301 - 2.699)/(2 - 0) = 0.301$ per unit of $t$}
  \item $b = 10^{0.301} \approx $ \ans{$2$} \quad and \quad $a = 10^{2.699} \approx $ \ans{$500$}
  \item Exponential model: $P(t) = $ \ans{$500 \cdot 2^t$}
  \item Verify: $P(2) = $ \ans{$500 \cdot 4 = 2000$}. Does this match? \ans{Yes}
\end{enumerate}

\item \textbf{Comparison.}  What advantage does Scientist B (semi-log plot) have?

\ansline{Scientist B's approach is more direct and does not require guessing an appropriate
constant to add to the $y$-values. According to EK 2.15.A.2, a constant never needs to be
added to the dependent variable values on a semi-log plot to reveal that an exponential model
is appropriate. If $\log(y)$ vs.\ $x$ is linear, an exponential model fits --- no constant
required.}

\item \textbf{AP-style multiple choice.}  For $y = 8 \cdot 6^x$, slope of the linearized
      semi-log model:

\vspace{6pt}
\begin{enumerate}[leftmargin=2em, itemsep=4pt, label=(\Alph*)]
  \item $\log_{10}(8)$
  \item \ans{$\log_{10}(6)$} \quad \textbf{(B) --- slope $= \log_{10}(b) = \log_{10}(6)$}
  \item $\log_{10}(48)$
  \item $6$
\end{enumerate}

\end{enumerate}

\vspace{0.1in}

\begin{tcolorbox}[colback=goldbg, colframe=goldacc, arc=2mm, left=3mm, right=3mm,
    top=2mm, bottom=2mm, title=\textbf{Extension}, fonttitle=\bfseries]
Natural-log linearization of $y = ab^x$:

\ansline{$\ln(y) = \ln(a) + x\ln(b)$. Slope $= \ln(b)$ and intercept $= \ln(a)$.
By change of base: $\ln(b) = \log_{10}(b)/\log_{10}(e) \approx \log_{10}(b)/0.4343$.
So the slope using $\ln$ is larger than with $\log_{10}$ by a factor of $1/\log_{10}(e)
\approx 2.303$. Similarly the intercept is scaled. The shape is the same; only the scale
of the axes differs.}
\end{tcolorbox}

\end{document}
